% ============================================
% DOC-25
% Extraído de main.pandoc.tex
% ============================================
\section*{\textbf{DOC-25} --- O Homem do Pano Preto}
\markboth{DOC-25 --- O Homem do Pano Preto}{}

\noindent \textit{Documento extraído do arquivo pericial.}

\vspace{0.5em}\noindent\rule{\textwidth}{0.5pt}\vspace{0.5em}

\noindent {\footnotesize\texttt{ARQUIVO P\'UBLICO $\cdot$ PROVA DOCUMENTAL $\cdot$ PROCESSO 1.260}}\\[-0.2em]

\noindent \textbf{Arquivo Público Mineiro --- Depoimento Informal}\\
\noindent \textbf{Depoente:} Wellington Alves da Silva, 34 anos, segurança patrimonial (terceirizado).\\
\noindent \textbf{Entrevistador:} Dra. Lúcia Mendes (gravado com autorização).\\
\noindent\textbf{Data:} 15 de janeiro de 2026. \textbf{Local:} Sala de descanso dos funcionários, APM.

\noindent \textbf{Lúcia:} O senhor estava na recepção quando o objeto foi devolvido?

\noindent \textbf{Wellington:} Tava. Segunda-feira, 5 de janeiro. Umas 18h30. A recepção já ia fechar. Eu tava fazendo a ronda final.

\noindent \textbf{Lúcia:} Viu quem entregou?

\noindent \textbf{Wellington:} Vi. Um homem. Idade difícil de dizer. Podia ter 50. Podia ter 70. O rosto era... como explicar... era um rosto que a gente esquece enquanto olha. Eu olhei pra ele e um segundo depois já não lembrava o nariz, a boca, nada. Só os olhos. Os olhos eram a única coisa que ficava.

\noindent \textbf{Lúcia:} Que cor?

\noindent \textbf{Wellington:} (pausa longa) O senhor já viu olho amarelo? Não olho de doente. Olho de bicho. De bicho que enxerga no escuro. Era assim. Amarelo claro, quase dourado, com a pupila preta no meio. A pupila não era redonda. Era... (gestícula vertical) ...fendida. Igual cabra.

\noindent \textbf{Lúcia:} Ele falou alguma coisa?

\noindent \textbf{Wellington:} Falou. Disse: \textquotedblleft{}Isto é para a doutora.\textquotedblright{} Só isso. Colocou o pacote no balcão, virou e saiu. O pacote era um pano preto. Daqueles panos grossos, sabe? De cobrir santo. Quando eu toquei, tava morno. Como se tivesse sido segurado perto do corpo por muito tempo.

\noindent \textbf{Lúcia:} O senhor abriu?

\noindent \textbf{Wellington:} Deus me livre. Eu vi o que era pelo formato. Um caderno. Um caderno velho. Mas não foi o caderno que me deu medo. Foi o pano. Quando eu fui dobrar o pano pra guardar, ele tava manchado. Não era mancha de tinta. Era um negócio escuro, meio marrom, com um cheiro... (pausa) ...a senhora vai achar que eu tô louco.

\noindent \textbf{Lúcia:} Pode falar.

\noindent \textbf{Wellington:} O pano tinha cheiro de mato molhado. Mas não mato de jardim. Mato de lugar onde enterraram coisa. Aquele cheiro de terra virada. Quando eu respirei, senti gosto de ferrugem na língua.

\noindent \textbf{Lúcia:} O senhor guardou o pano?

\noindent \textbf{Wellington:} Joguei fora. Na mesma noite. Dentro de um saco amarrado. (pausa) No dia seguinte, o saco tinha sumido da lixeira.

\noindent \textbf{Lúcia:} Sumido?

\noindent \textbf{Wellington:} A lixeira tava vazia. O saco não. Só o pano. O pano saiu de dentro do saco amarrado e sumiu. A senhora pode perguntar pro Manoel, da limpeza. Ele viu. Disse que nunca aconteceu uma coisa dessas. Eu falei: \textquotedblleft{}Agora aconteceu.\textquotedblright{}

\noindent \textit{Nota do arquivista: Wellington Alves da Silva pediu demissão 10 dias após este depoimento. O formulário de desligamento não registra motivo. Seu último dia de trabalho foi 25 de janeiro de 2026 --- 13 dias antes de o diário ser retirado do arquivo pela segunda vez.}
\noindent \textbf{Adendo --- Identificação posterior (Dra. Lúcia Mendes, 4 de fevereiro de 2026):}

\noindent Cruzei a descrição de Wellington com a fotografia do crachá do Dr. Heitor Oliveira (1979) e com a foto encontrada dentro do diário --- um homem mais velho, sentado num banco de praça, com um livro no colo.

\noindent A estrutura óssea é a mesma. O formato do crânio, o ângulo da mandíbula, a distância entre os olhos. O homem do pano preto \textbf{é} Heitor Oliveira.

\noindent Ele teria 81 anos em janeiro de 2026. A idade corresponde.

\noindent Mas a pupila fendida descrita por Wellington --- olho de cabra, vertical --- não é humana. Não é envelhecimento. É transformação. A entidade não o matou. Ela o \textit{manteve}. 47 anos de hospedagem mudaram o corpo dele. O olho direito já não é um olho humano. É uma fenda. Um portal. O diário usou Oliveira como entregador --- porque ele era o único que ainda podia andar entre os vivos sem levantar suspeitas.

\noindent O paciente 1.261 não desapareceu. Ele foi \textit{guardado} para o momento da entrega.

\noindent \textit{Nota final: Dra. Lúcia Mendes anexou esta identificação ao processo em 4 de fevereiro de 2026. Sete meses depois, ela desapareceu. A última pessoa a vê-la foi um segurança de rodoviária, que descreveu uma mulher com olho direito amarelo embarcando num ônibus para o Ceará.}

\vspace{0.5em}
\noindent\textit{\textbf{↪ Rotas:} \rota{DOC-01} \rotasep \rota{LUC-02} \rotasep \rota{DOC-26}}
